\begin{frame}{Änderungsanträge (EMPL)}
\tiny
\vspace{-0.45cm}
\renewcommand{\arraystretch}{0.9}
\setlength{\tabcolsep}{2pt}
\setlength{\abovetopsep}{0pt}
\setlength{\belowrulesep}{0pt}
\begin{tabularx}{\textwidth}{@{}>{\centering\arraybackslash}m{1.55cm} >{\RaggedRight\arraybackslash}m{\dimexpr\linewidth-1.55cm-2\tabcolsep\relax}@{}}
\toprule
\multicolumn{1}{m{1.55cm}}{\centering\textbf{Fraktion}} & \multicolumn{1}{m{\dimexpr\linewidth-1.55cm-2\tabcolsep\relax}}{\centering\textbf{Änderungsantrag}} \\
\midrule
\includegraphics[width=2cm,height=1cm,keepaspectratio]{Bilder/EVP.png} & a) der frühzeitige Zugang zum Arbeitsmarkt sowie die Bereitstellung \textcolor{green}{\uline{von Qualifizierungsangeboten sowie der verpflichtende Besuch }}von Sprachkursen \textcolor{red}{\sout{und Qualifizierungsangeboten}}\textcolor{green}{\uline{mit darauffolgenden verpflichtenden Prüfungen der erworbenen Sprachkenntnissen}}, \\
\includegraphics[width=2cm,height=1cm,keepaspectratio]{Bilder/SD.png} & a) der frühzeitige Zugang \textcolor{green}{\uline{)}}zum Arbeitsmarkt sowie die Bereitstellung von Sprachkursen und Qualifizierungsangeboten\textcolor{green}{\uline{. Der Vorgang zur Arbeitsfähigkeitsprüfung soll nach spätestens zwei Jahren vorankommen.}}, \\
\includegraphics[width=2cm,height=1cm,keepaspectratio]{Bilder/RE.png} & d) die Gewährleistung \textcolor{red}{\sout{von sozialen Dienstleistungen und }}der \textcolor{green}{\uline{notwendigen }}Gesundheitsversorgung, \\
\includegraphics[width=2cm,height=1cm,keepaspectratio]{Bilder/PfE.png} & 1. Die Mitgliedstaaten \textcolor{red}{\sout{ergreifen}}\textcolor{green}{\uline{können auf freiwilliger Basis}} Maßnahmen zur Förderung der sozialen und wirtschaftlichen Integration von Schutzsuchenden und anerkannten Schutzberechtigten\textcolor{red}{\sout{, um deren aktive Teilhabe am gesellschaftlichen Leben sicherzustellen}}\textcolor{green}{\uline{ ergreifen}}. \\
\includegraphics[width=2cm,height=1cm,keepaspectratio]{Bilder/Green.png} & 3. Die Mitgliedstaaten berichten der Kommission jährlich über ihre Integrationsmaßnahmen.\textcolor{green}{\uline{ Ein Prüfungsverfahren soll eingeführt werden, welches prüft, in wie fern die verschiedenen Mitgliedsstaaten sich gegenseitig solidarisch unterstützen, wenn es um entstehnde Last und Arbeit durch Asylanträge geht. Zusätzlich soll ein Solidaritäsmechanismus und weitere Unterstützung umgesetzt werden, zum Beispiel durch Gelder der Mitgliedsstaaten im Sinne der Integration und für die Schutzsuchenden.}} \\
\bottomrule
\end{tabularx}
\end{frame}

